\documentclass[11pt,a4paper]{article}

\usepackage[utf8]{inputenc}
\usepackage[T1]{fontenc}
\usepackage[ngerman]{babel}
\usepackage{amsmath, amssymb}
\usepackage{graphicx}
\usepackage{booktabs}
\usepackage[margin=2.5cm]{geometry}
\usepackage{hyperref}

\title{Neuronensimulation mit dem Hodgkin-Huxley-Modell}
\author{Schlomi \and Projektpartner}
\date{\today}

\begin{document}
\maketitle

% Vorgabe: kurzer Bericht, max. 3-4 Seiten ohne Figuren.
% Diskussion der Problemstellung und der Hauptresultate.

\section{Einleitung}
% Kurz: Problemstellung und Ziel des Projekts.

\section{Hodgkin-Huxley-Modell}

\subsection{Grundlegende Überlegungen (Aufgabe 1)}
% 1a) Herleitung von (11)-(13) aus (8)-(10):
% mit tau = 1/(alpha+beta) und x_inf = alpha/(alpha+beta).
% 1b) Sinnvolle Anfangsbedingungen fuer n, m, h (inaktives Neuron).

\subsection{Numerische Methoden (Aufgabe 2)}
% Euler vs. Runge-Kutta (RK4) vs. scipy.odeint.
% Stabilitaet, Schrittweite, Performance.
% Stromvariation und Stromimpuls; Rolle der Gatingvariablen.

\section{Neuronales Netz von Grund auf (Aufgabe 3)}
% Aufbau, Wichtungsfaktoren, Lernalgorithmus, Robustheit.

\section{Neuronale Netze mit Keras (Aufgaben 4 \& 5)}
% Schachbrett mit Keras; MNIST/fashion-MNIST; Confusion Matrix; Hidden-Layer.

\section{Fehlerabschätzung und Diskussion}
% Genauigkeit/Zuverlaessigkeit der Rechnungen; analytische Grenzfaelle als Test.

\section{Fazit}

\end{document}
